\subsection{Selecting enhancers for multi-cell expression programs}
\label{app:multicell-programming}

\paragraph[Assay, candidates and source split.]{\phantomsection\label{app:multicell-split}Assay, candidates and source split.}
We use Atlas reporter measurements \citep{castillohair2025atlas}
(\href{https://zenodo.org/records/17410822}{Zenodo record 17410822}).
The sixteen requests are its eight double-target and eight triple-target
Fast SeqProp (FSP) programs, without outcome filtering. Every request receives
all 2,411 distinct 145-nt sequences from the double/triple FSP and genomic
categories, regardless of original target. Exact DNA duplicates with identical
counts are collapsed. The twelve cell lines are NT2\_D1,
GM12878, 786\_O, SKNSH, WERI\_Rb1, SJCRH30, HepG2, K562, MCF7, HeLaS3, HEK293
and HMC3; HEK293 denotes the authors' HEK293T key. Retina is excluded because
it has only one RNA measurement.

Calibration uses single-target FSP, Deep Exploration Network, genomic and
tunable-FSP constructs. Families are connected components linking shared
background/DHS identifiers and sequences within global Levenshtein distance 29
on either strand. Excluding
families that touch candidates leaves 5,246 source sequences. A deterministic
family-hash split assigns 4,223 to training and 1,023 to validation; refitting
uses all 5,246.

\paragraph{Measured activity.}
For each cell and repeat $r\in\{1,2\}$, activity is
\[
 y_{icr}=\log_2\frac{R_{icr}+0.5}{D_{icr}+0.5}
 -\operatorname{median}_{j\in\mathcal N}
   \log_2\frac{R_{jcr}+0.5}{D_{jcr}+0.5},
\]
using UMI-deduplicated counts and the published 35 negative controls $\mathcal N$,
which had undergone pooled outlier removal from forty. DNA preparations follow
Supplementary Table~8. We retain the source DNA QC ($>20$ total reads in every
preparation) without an additional RNA outcome filter. The two RNA measurements come from separate
transfections; DNA repeats are technical sequencing replicates of shared plasmid
preparations. Normalization is separate for each cell and repeat.

\paragraph{Model predictions.}
Inserts in the published orientation receive the 16-nt left and 6-nt right cloning
junctions, with Ns outside. Inputs are 16,384 bases for AlphaGenome, 196,608 for
Enformer and 1,024 for NTv3 post; they represent short inserts, not complete
reporters or genomic loci. We use all 4,418 AlphaGenome human ATAC, DNase, CAGE,
RNA-seq and histone/TF ChIP tracks, 5,313 Enformer human outputs and 7,362 NTv3
human BigWig outputs. Released nonlinear heads precede insert averaging, with
128-bp bins weighted by overlap. NTv3 uses its official 384-position crop.
Pooled outputs receive $\log(1+x)$, without reverse-complement ensembling or
projection. These native outputs require assay calibration.
AlphaGenome hidden-feature readouts use the same calibration and candidate catalogue.

\paragraph{Calibration and controls.}
Source targets average both normalized repeats. Training statistics standardize
features; dimensions with standard deviation at most $10^{-7}$ are removed,
and features are divided by the square root of retained width. Ridge minimizes
sample-mean summed squared error plus $\lambda\lVert W\rVert_F^2$, leaving the
intercept unpenalized. Mean twelve-cell validation MSE selects among seventeen
penalties $10^{-8},10^{-7.5},\ldots,1$; scaling and coefficients are then refitted
on all source rows. The same procedure calibrates the authors' 64 DHS predictions
and strand-invariant 1--6-mer frequencies.

The CNN has three convolutional layers and twelve outputs. Source validation
selects the learning rate and stopping epoch for each of three seeds. Models
are then refitted on all source sequences and their predictions averaged.
Candidate measurements are used only for evaluation; architecture and training
settings are supplied with the code.

\paragraph{Selection and evaluation.}
Predictors rank candidates by Equation~\ref{eq:multicell-utility}, selecting ten
with ascending-ID ties and no abstention.
The same nominations are scored in both repeat measurements, averaged within program
and then equally across sixteen programs.
Controls are each predictor's minimum-ON-only ranking, uniform catalogue selection
and uniform selection within the author-intended target cohort; random utilities
are exact expectations. Experimental selection ranks with one repeat measurement and
scores the other, reciprocally; these evaluations are dependent.
Figure~\ref{fig:frontier-programming}A subtracts
intent-matched random utility per program/repeat; markers show repeats, not
confidence intervals.

\begin{figure}[!htbp]
 \centering\capfigure{GenomicDesign}
 \includegraphics[width=.92\linewidth]{multicell_selector_control.pdf}
 \capcaption{GenomicDesign}{\textbf{Considering unwanted activity improves sequence selection.}
 The same full-profile utility evaluates target-only and full-profile selections,
 averaged over ten nominations, sixteen requests and both repeat measurements. Open
 markers rank by the weakest requested activity alone; filled markers also
 penalize the strongest unwanted activity. The experimental reference selects
 on the other repeat measurement, reciprocally. Connecting lines are policy differences,
 not uncertainty intervals.}
 \label{fig:atlas-selector-control}
\end{figure}

Figure~\ref{fig:frontier-programming}B shows the ten AlphaGenome nominations for
HepG2/K562 in predicted order, with profiles averaged across both repeats. This comparatively successful request illustrates the setting, not aggregate performance.
Figure~\ref{fig:atlas-selector-control} compares full-profile and ON-only selection.
